\cvsection{Presentations (\cvstar\ denotes presenter)}

\cvsubsection{Invited Oral Presentations}
\begin{cvnumbered}
  \item Models of infectious disease dynamics: Making predictions to support public health response. \textbf{E Howerton\cvstar}. \textit{Department of Biology and Environmental Data Science Seminar, James Madison University}. April 11, 2025.

  \item Multi-model scenario projections to support public health in real time: Challenges and opportunities from the U.S. Scenario Modeling Hubs. \textbf{E Howerton\cvstar}. \textit{Disease Decision Analysis \& Research Group, U.S. Geological Survey}. February 10, 2025.

  \item Informing public health response in the face of uncertainty: Lessons from the U.S. Scenario Modeling Hub. \textbf{E Howerton\cvstar}, C Viboud, and J Lessler on behalf of the COVID-19 Scenario Modeling Hub. \textit{Isaac Newton Institute, Modelling and Inference for Pandemic Preparedness---A Focused Workshop}. August 5--9, 2024. \href{https://www.youtube.com/watch?v=BSKv0QbDhoc}{Recording}.

  \item Synergistic interventions to control COVID-19. \textbf{E Howerton\cvstar}, MJ Ferrari, ON Bjørnstad, TL Bogich, RK Borchering, CP Jewell, JD Nichols, WJM Probert, MJ Tildesley, MC Runge, C Viboud, and K Shea. \textit{Joint Mathematics Meetings}. April 6--8, 2022.

  \item Synergistic interventions to control COVID-19. \textbf{E Howerton\cvstar}, MJ Ferrari, ON Bjørnstad, TL Bogich, RK Borchering, CP Jewell, JD Nichols, WJM Probert, MJ Tildesley, MC Runge, C Viboud, and K Shea. \textit{Department of Biology, Pennsylvania State University}. Invited presentation by the Jeanette Ritter Mohnkern Award winner. August 24, 2021.
\end{cvnumbered}

\cvsubsection{Oral Presentations}
\begin{cvnumbered}
  \item The honeymoon is over: Dynamic impacts of abrupt reductions in vaccine uptake. \textbf{E Howerton\cvstar}, BT Grenfell, CJE Metcalf. \textit{Epidemics\textsuperscript{10}}. December 1, 2025.

  \item Community dynamics of human respiratory pathogens: The HMPV story. \textbf{E Howerton\cvstar}. \textit{Department of Ecology and Evolutionary Biology, Princeton University}. February 7, 2025.

  \item Modeling the impact of RSV interventions on hMPV burden and dynamics. \textbf{E Howerton\cvstar}. \textit{Fogarty International Center EPS Modeling Seminar, National Institutes of Health}. October 16, 2024.

  \item An introduction to human metapneumovirus. \textbf{E Howerton\cvstar}. \textit{Life Course Immunity: A Multi-Pathogen Perspective}. February 9, 2024.

  \item Informing pandemic response in the face of uncertainty: An evaluation of the U.S. COVID-19 Scenario Modeling Hub. \textbf{E Howerton\cvstar}, C Viboud, and J Lessler on behalf of the COVID-19 Scenario Modeling Hub. \textit{Epidemics\textsuperscript{9}}. November 28, 2023.

  \item What does it mean to evaluate a scenario projection? Lessons from the COVID-19 Scenario Modeling Hub. \textbf{E Howerton\cvstar}, J Lessler, and C Viboud on behalf of the COVID-19 Scenario Modeling Hub contributors. \textit{Royal Society Meeting on Forecasting Infectious Disease Incidence}. March 15, 2023.

  \item Collaborative open hubs for infectious disease modeling. C Viboud\cvstar, \textbf{E Howerton}, A Vespignani, and J Lessler on behalf of the COVID-19 Scenario Modeling Hub and Ebola Forecast Challenge contributors. \textit{Royal Society Meeting on Forecasting Natural and Social Systems}. March 14, 2023.

  \item Optimizing management decisions for control of infectious disease outbreaks. \textbf{E Howerton\cvstar}. \textit{Northeastern University}. February 6, 2023.

  \item A look back at fifteen rounds of SMH: Evaluation. \textbf{E Howerton\cvstar}. \textit{COVID-19 Scenario Modeling Hub Meeting 2022---Past, Present, and Future of Scenario Modeling for the U.S.} September 14, 2022.

  \item Looking back at one year of the COVID-19 Scenario Modeling Hub. \textbf{E Howerton\cvstar}, on behalf of the COVID-19 Scenario Modeling Hub Team. \textit{MIDAS Annual Meeting}. September 8, 2022.

  \item Misapplied management makes matters worse: Spatially explicit control strategies that leverage biotic interactions may more effectively prevent invader spread. \textbf{E Howerton\cvstar}, T Langkilde, and K Shea. August 17, 2022.

  \item How to get the best out of multiple disease models. \textbf{E Howerton\cvstar}, MC Runge, TL Bogich, RK Borchering, H Inamine, J Lessler, LC Mullany, WJM Probert, CP Smith, S Truelove, C Viboud, and K Shea. \textit{Ecology and Evolution of Infectious Diseases}. June 6--9, 2022. \href{https://www.youtube.com/watch?v=zLB9zyG-koI}{Recording}.

  \item Scenario Modeling Hub round 13 preview and evaluation. J Lessler\cvstar\ and \textbf{E Howerton\cvstar} on behalf of the COVID-19 Scenario Modeling Hub. \textit{U.S. Centers for Disease Control and Prevention Council of State and Territorial Epidemiologists}. April 5, 2022.

  \item Comparing aggregation methods for the COVID-19 Scenario Modeling Hub. \textbf{E Howerton\cvstar}. \textit{COVID-19 Scenario Modeling Hub Weekly Meeting}. June 18, 2021.
\end{cvnumbered}

\cvsubsection{Poster Presentations}
\begin{cvnumbered}
  \item Dynamic consequences of declining vaccination coverage in a heterogeneous world. \textbf{E Howerton\cvstar}, A Mahmud, DJD Earn, BT Grenfell, CJE Metcalf. \textit{Ecology and Evolution of Infectious Disease Dynamics}. June 3, 2026.

  \item Cross-protection and the lags between RSV and human metapneumovirus outbreaks. \textbf{E Howerton\cvstar}, TC Williams, J-S Casalegno, CJE Metcalf, SW Park, BT Grenfell. \textit{Ecology and Evolution of Infectious Disease Dynamics}. June 25, 2024.

  \item Informing pandemic response in the face of uncertainty: An evaluation of the U.S. COVID-19 Scenario Modeling Hub. \textbf{E Howerton\cvstar}, K Shea, C Viboud, and J Lessler on behalf of the COVID-19 Scenario Modeling Hub. \textit{Ecology and Evolution of Infectious Disease Dynamics}. May 23, 2023.

  \item Theory of combining probabilistic projections with applications in epidemiology. \textbf{E Howerton\cvstar}, MC Runge, TL Bogich, RK Borchering, H Inamine, J Lessler, LC Mullany, WJM Probert, CP Smith, S Truelove, C Viboud, and K Shea. \textit{Epidemics\textsuperscript{8}}. November 29--December 1, 2021.

  \item {\addfontfeatures{LetterSpace=-0.5}Synergistic interventions to control COVID-19: Mass testing and isolation mitigates reliance on distancing. \textbf{E Howerton\cvstar}, MJ Ferrari, ON Bjørnstad, TL Bogich, RK Borchering, CP Jewell, JD Nichols, WJM Probert, MJ Tildesley, MC Runge, C Viboud, and K Shea. \textit{Ecology and Evolution of Infectious Diseases}. June 14--17, 2021.}
\end{cvnumbered}
